% C7: Statistical Significance Tests
% Formal hypothesis tests for banditGPT vs baselines.
% Cross-references: Figure 4, C4 (detailed Pareto results)

\subsection{Statistical Significance Tests}
\label{appendix:statistical_tests}

This subsection provides formal hypothesis tests for the main Figure~4 claim
that banditGPT-Hybrid outperforms RouteLLM-MF on the holdout set.

\subsubsection{Test Design}

\paragraph{Challenge.}
banditGPT is stochastic (seed-dependent exploration), producing 20
independent trial means.  RouteLLM-MF is deterministic at each threshold,
producing a single aggregate reward.  A paired per-prompt test between the
two is ill-defined because RouteLLM's per-prompt selections are not available
(we use only its published aggregate rewards from the threshold sweep).

\paragraph{Approach.}
We use a \textbf{one-sample $t$-test} to test whether banditGPT's mean reward
(across 20 seeds) is significantly different from RouteLLM-MF's best
aggregate reward ($0.8827$ at the Pareto-optimal threshold).  We supplement
with a Wilcoxon signed-rank test (non-parametric) and a bootstrap CI.

For the \textbf{Static-Mixtral} baseline, we compute a per-prompt paired
$t$-test using the seed-averaged per-prompt banditGPT reward against
Mixtral's known per-prompt reward.

\subsubsection{Results}

\begin{table}[h]
\centering
\caption{Statistical significance tests ($\lambda{=}0$, 750-prompt holdout, 20 seeds).}
\label{tab:significance}
\small
\begin{tabular}{@{}lccccc@{}}
\toprule
\textbf{Comparison} & \textbf{$\Delta$ Reward} & \textbf{$t$-stat} & \textbf{$p$-value} & \textbf{Cohen's $d$} & \textbf{Wilcoxon $p$} \\
\midrule
banditGPT vs.\ RouteLLM-MF & $+0.019$ & 7.34 & $6.0 \times 10^{-7}$ & 1.64 & $2.2 \times 10^{-4}$ \\
banditGPT vs.\ Static Mixtral & $+0.079$ & 6.52 & $1.3 \times 10^{-10}$ & 0.24 & $4.0 \times 10^{-9}$ \\
\bottomrule
\end{tabular}
\end{table}

\paragraph{Bootstrap CI.}
The 95\% bootstrap CI for banditGPT's mean reward is $[0.896, 0.906]$,
which entirely excludes RouteLLM-MF's best point ($0.883$).

\paragraph{Oracle gap closure.}
banditGPT closes $\frac{0.902 - 0.823}{0.953 - 0.823} = 60.3\%$ of the
Oracle--Mixtral gap, vs.\ RouteLLM-MF's 45.9\%.

\subsubsection{Interpretation}

\paragraph{1.\ Large, consistent advantage over RouteLLM-MF.}
Cohen's $d = 1.64$ is well above the ``large'' threshold ($d \geq 0.8$),
and 19 of banditGPT's 20 trial means (range: 0.865--0.913) exceed
RouteLLM-MF's best threshold reward (0.883).  The 95\% bootstrap CI
$[0.896, 0.906]$ entirely excludes the RouteLLM-MF reference, providing
strong evidence that banditGPT's advantage is not due to favorable
seed selection.

\paragraph{2.\ Targeted per-prompt improvement over Mixtral.}
The per-prompt effect against static Mixtral (Cohen's $d = 0.28$) is
``small'' by convention, which reflects the nature of LLM routing: most
prompts yield similar reward under either model (ties dominate).
banditGPT's aggregate advantage comes from correctly routing the minority
of prompts with clear model preferences --- precisely the prompts where
contextual routing adds value.

\paragraph{3.\ Conservative interpretation.}
The one-sample $t$-test against RouteLLM-MF treats the 0.883 reference as
exact.  In reality, RouteLLM's aggregate reward has some uncertainty from
the threshold-selection procedure (which threshold is ``best'' depends on
the holdout set), which inflates the $t$-statistic and Cohen's $d$.
We therefore emphasize the bootstrap CI as the most conservative measure:
even the lower bound (0.896) exceeds the RouteLLM-MF reference by 1.3
percentage points, confirming significance under any reasonable
uncertainty in the baseline.

\paragraph{4.\ Scope: quality-focused regime only.}
These tests compare methods at $\lambda{=}0$ (pure quality optimization).
The Pareto frontier analysis (Figure~4; Appendix~\ref{appendix:pareto_detailed})
reveals a more nuanced picture across the full cost--quality plane: at low
budgets (cost $< \$0.005$), RouteLLM-MF's prompt-level discrimination
makes it competitive with banditGPT, while the large, significant advantage
reported here applies to the moderate-to-high budget regime where banditGPT's
online aggregate calibration dominates.

\subsubsection{Reproducibility}

Code: \texttt{experiments/04\_figure/run\_statistical\_tests.py}

\begin{itemize}[noitemsep]
    \item Runtime: ${\sim}33$ seconds (Apple M2)
    \item Results: \texttt{results/statistical\_tests.json}
\end{itemize}
